\documentclass[11pt]{article}
\usepackage[margin=1.1in]{geometry}
\usepackage{amsmath,amssymb,amsthm,mathtools}
\usepackage{hyperref}

\theoremstyle{plain}
\newtheorem{theorem}{Theorem}
\newtheorem{proposition}[theorem]{Proposition}
\newtheorem{lemma}[theorem]{Lemma}
\newtheorem{corollary}[theorem]{Corollary}
\theoremstyle{definition}
\newtheorem{definition}[theorem]{Definition}
\newtheorem{remark}[theorem]{Remark}
\newtheorem{example}[theorem]{Example}

\DeclareMathOperator{\ord}{ord}
\DeclareMathOperator{\hgt}{ht}
\DeclareMathOperator{\ad}{ad}
\newcommand{\LL}{\Lambda}
\newcommand{\eps}{\varepsilon}
\newcommand{\Q}{\mathbb{Q}}
\newcommand{\Z}{\mathbb{Z}}
\newcommand{\ind}[1]{[\![#1]\!]}

\title{The ribbon operator is not a differential operator\\[2pt]
\large $\ord(N_e)=\infty$, and the order filtration is discontinuous at $t=-1$}
\author{Clio}
\date{6 September 2026}

\begin{document}
\maketitle

\begin{abstract}
Let $R_e(t)$ be the operator on symmetric functions that adds a connected border
strip of size $e$ with weight $t^{\hgt}$, let $g_e(t)=M[P_{(e)}(X;-t)]$ be
multiplication by the one-row Hall--Littlewood polynomial, and let
$E_e=g_e-R_e$ be the disconnected complement.  Write $N_e:=\partial_t E_e|_{t=-1}$,
the operator that adds a size-$e$ broken border strip with \emph{exactly two}
components, signed by height.  We compute two closed forms for the iterated
commutator of $R_e(t)$ with multiplication by $p_1$, evaluated on two hook
families, and deduce:
\[
  \ord\bigl(R_e(c)\bigr)=\infty\ \text{ for every scalar }c\neq-1,
  \qquad \ord\bigl(R_e(-1)\bigr)=\ord(M_{p_e})=0,
\]
and $\ord(N_e)=\ord(E_e)=\infty$ for every $e\ge2$.  So $N_e$ is not a derivation,
is not of any finite order, and is not a differential operator on $\LL$ at all.
Consequently none of $R_e(t)$, $E_e(t)$, $N_e$ lies in the algebra generated by
multiplication operators $M_f$ and skewing operators $f^{\perp}$.  The mechanism
is a single exceptional chain --- the one that adds the strip to the empty
partition --- and the factor $(1+t)$ in the first closed form is exactly the
two ways that chain can sit.
\end{abstract}

\section{Setting and conventions}

Fix the ground field $K=\Q(t)$ and let $\LL=K[p_1,p_2,\dots]$ be the ring of
symmetric functions, with Schur basis $\{s_\lambda\}$.  For $f\in\LL$ write
$M_f$ for multiplication by $f$ and $f^{\perp}$ for its Hall-inner-product
adjoint.

\begin{definition}[order filtration]
$\mathcal D^{\le-1}=0$, and for $d\ge0$,
\[
  \mathcal D^{\le d}
  =\bigl\{X\in\operatorname{End}_K(\LL)\ :\ [X,M_f]\in\mathcal D^{\le d-1}
    \text{ for all }f\in\LL\bigr\}.
\]
Put $\ord(X)=\min\{d:X\in\mathcal D^{\le d}\}$, and $\ord(X)=\infty$ if no such
$d$ exists.
\end{definition}

This is Grothendieck's filtration by differential order for the commutative ring
$\LL$; under $\LL=K[p_1,p_2,\dots]$ it is the order of $X$ as a differential
operator in the power sums.  We use three standard facts.

\begin{lemma}\label{lem:std}
\begin{enumerate}
\item[(i)] $\ord(M_f)=0$ for every $f\in\LL$.
\item[(ii)] $\mathcal D^{\le a}\mathcal D^{\le b}\subseteq\mathcal D^{\le a+b}$;
in particular $\ord(XY)\le\ord(X)+\ord(Y)$, and $\mathcal D^{\le d}$ is a
two-sided $\mathcal D^{\le0}$-module.
\item[(iii)] $\ord(p_\rho^{\perp})=\ell(\rho)$, and more generally
$\ord(f^{\perp})=\deg_p f:=\max\{\ell(\rho):[p_\rho]f\neq0\}$
\emph{(Macdonald I.5)}.
\end{enumerate}
\end{lemma}

\begin{lemma}[generators suffice]\label{lem:gen}
$X\in\mathcal D^{\le d}$ if and only if
$[\,\cdots[[X,M_{p_{a_1}}],M_{p_{a_2}}],\dots,M_{p_{a_{d+1}}}]=0$
for all $a_1,\dots,a_{d+1}\ge1$.
\end{lemma}

\begin{proof}
Necessity is immediate.  For sufficiency, induct on $d$.  The case $d=0$ says
$[X,M_{p_a}]=0$ for all $a$; since $[X,M_{fg}]=[X,M_f]M_g+M_f[X,M_g]$ and
$[X,M_c]=0$ for $c\in K$, an induction on the number of $p$-factors of a
monomial $p_\rho$ gives $[X,M_{p_\rho}]=0$, hence $[X,M_f]=0$ for all $f$ by
linearity, i.e.\ $X\in\mathcal D^{\le0}$.  For $d\ge1$ the same product rule,
together with Lemma~\ref{lem:std}(ii) (which makes $\mathcal D^{\le d-1}$ stable
under left and right multiplication by $M_g$), propagates the hypothesis from
generators to all $f$; the inductive hypothesis then places $[X,M_f]$ in
$\mathcal D^{\le d-1}$.
\end{proof}

Throughout we write
\[
  \ad(X):=[X,M_{p_1}],\qquad
  \ad^{\,d}(X)=\sum_{r=0}^{d}(-1)^r\binom{d}{r}\,M_{p_1}^{\,r}\,X\,M_{p_1}^{\,d-r},
\]
the second equality because the $M_{p_1}$ commute with each other.  By
Lemma~\ref{lem:gen}, \emph{if $\ad^{\,d}(X)\neq0$ for every $d$ then
$\ord(X)=\infty$}; this is the only direction we use.

\subsection{The operators}

Let $\lambda\subseteq\mu$ with $|\mu/\lambda|=e$.  Call $\mu/\lambda$ a
\emph{broken border strip} if it contains no $2\times2$ square; let
$m(\mu/\lambda)$ be its number of edge-connected components and
$\hgt(\mu/\lambda)=\sum_{\text{components }c}(\#\text{rows}(c)-1)$.  Define
\begin{align*}
  R_e(t)\,s_\lambda &=\textstyle\sum_{m(\mu/\lambda)=1} t^{\hgt(\mu/\lambda)}\,s_\mu,\\
  g_e(t)\,s_\lambda &=\textstyle\sum_{\mu/\lambda\ \mathrm{broken}}
        (1+t)^{m(\mu/\lambda)-1}\,t^{\hgt(\mu/\lambda)}\,s_\mu,
  \qquad E_e:=g_e-R_e .
\end{align*}
At $t=-1$ the Murnaghan--Nakayama rule gives $R_e(-1)=M_{p_e}$ and $E_e(-1)=0$.

\begin{proposition}[the dictionary, Konvalinka; Halverson--Ram]\label{prop:hl}
$g_e(t)=M\bigl[P_{(e)}(X;-t)\bigr]$, multiplication by the one-row
Hall--Littlewood polynomial at parameter $q=-t$.  Explicitly, with
$[n]_{q}=1+q+\dots+q^{n-1}$,
\begin{equation}\label{eq:Prow}
  P_{(e)}(X;-t)\;=\;\sum_{\rho\vdash e}\frac{(1+t)^{\ell(\rho)-1}}{z_\rho}
  \Bigl(\prod_i [\rho_i]_{-t}\Bigr)\,p_\rho .
\end{equation}
In particular $\ord(g_e)=0$.
\end{proposition}

\begin{proof}[Derivation of \eqref{eq:Prow}]
$P_{(r)}(X;q)=q_r(X;q)/(1-q)$ where
$\sum_{r\ge0}q_r(X;q)u^r=\prod_i\frac{1-qx_iu}{1-x_iu}
=\exp\bigl(\sum_{n\ge1}(1-q^n)p_nu^n/n\bigr)$.  Put $\eps:=1-q$, so
$1-q^n=\eps\,[n]_q$ and
$\sum_r q_ru^r=\exp\bigl(\eps B\bigr)$ with $B=\sum_n [n]_q p_nu^n/n$.  Hence
$P_{(e)}=\eps^{-1}[u^e]\sum_{j\ge1}\eps^jB^j/j!
=\sum_{j\ge1}\eps^{j-1}[u^e]B^j/j!$, and collecting compositions into
partitions gives \eqref{eq:Prow} with $q=-t$, $\eps=1+t$.
The identification $g_e=M[P_{(e)}(X;-t)]$ is the quantum Murnaghan--Nakayama
rule over broken ribbons.  $\ord(g_e)=0$ by Lemma~\ref{lem:std}(i).
\end{proof}

Set $\eps:=1+t$.  Since $t^h=(-1)^h(1-\eps)^h$,
\begin{equation}\label{eq:Rexp}
  R_e=\sum_{D\ge0}(-\eps)^D\,\Psi_{e,D},
  \qquad
  \Psi_{e,D}\,s_\lambda=\sum_{m(\mu/\lambda)=1}\binom{\hgt(\mu/\lambda)}{D}
  (-1)^{\hgt(\mu/\lambda)}s_\mu ,
\end{equation}
so $\Psi_{e,0}=M_{p_e}$ and $\Psi_{e,1}=-\partial_tR_e|_{t=-1}$.  Likewise
$E_e=\eps N_e+O(\eps^2)$ with
\begin{equation}\label{eq:Ndef}
  N_e\,s_\lambda:=\partial_t E_e|_{t=-1}\,s_\lambda
   =\sum_{m(\mu/\lambda)=2}(-1)^{\hgt(\mu/\lambda)}s_\mu ,
\end{equation}
because $(1+t)^{m-1}$ vanishes to order $m-1$ at $t=-1$: only $m=2$ survives the
first derivative.

\begin{remark}
The definition \eqref{eq:Ndef} is \emph{exactly two} components, not ``at least
two''.  This corrects the statement of $N_e$ in the session brief; the
open-questions file has it right.
\end{remark}

\subsection{The merging expansion}

Comparing $\eps$-Taylor coefficients in $g_e=R_e+E_e$ using
Proposition~\ref{prop:hl} and \eqref{eq:Rexp} gives, for each $D\ge0$, an
identity whose left side mixes component number with height and whose right
side is a multiplication operator.  We record the two we use.

\begin{proposition}\label{prop:merge}
For every $e\ge2$,
\[
  N_e-\Psi_{e,1}
  \;=\;M\Bigl[\ \partial_t P_{(e)}(X;-t)\big|_{t=-1}\Bigr]
  \;=\;M\Bigl[\ \tfrac12\sum_{a=1}^{e-1}p_a\,p_{e-a}\ -\ \tfrac{e-1}{2}\,p_e\ \Bigr],
\]
an operator of order $0$.  Hence $\ord(N_e)=\ord(\Psi_{e,1})$ whenever either is
positive, and $\ad^{\,d}(N_e)=\ad^{\,d}(\Psi_{e,1})$ for all $d\ge1$.
\end{proposition}

\begin{proof}
Since $\ord(g_e)=0$ and the $M_{p_a}$ do not depend on $t$, every $t$-Taylor
coefficient of $g_e$ is again a multiplication operator; the $\eps$-linear one
is $M[\partial_tP_{(e)}|_{t=-1}]$, and by \eqref{eq:Rexp}, \eqref{eq:Ndef} it
equals $-\Psi_{e,1}\cdot(-1)+N_e$, i.e.\ $N_e-\Psi_{e,1}$ up to the sign already
absorbed in \eqref{eq:Rexp}.  For the explicit form, expand \eqref{eq:Prow}:
$[n]_{-t}=n-\eps\binom{n}{2}+O(\eps^2)$, so writing
$\mathfrak m_\rho=\prod_i m_i(\rho)!$ and $z_\rho=\mathfrak m_\rho\prod_i\rho_i$,
the $\rho$-term of $P_{(e)}$ is
$\frac{\eps^{\ell(\rho)-1}}{\mathfrak m_\rho}
\bigl(1-\eps\frac{e-\ell(\rho)}{2}+O(\eps^2)\bigr)p_\rho$.  The $\eps^0$ part is
$p_e$; the $\eps^1$ part collects $\ell(\rho)=2$ at leading order and
$\rho=(e)$ at first order, giving
$\frac12\sum_{a=1}^{e-1}p_ap_{e-a}-\frac{e-1}{2}p_e$.
\end{proof}

\begin{remark}[why component number is only a filtration]
Proposition~\ref{prop:merge} is the precise form of the statement that
composing strip-additions merges components.  The $\eps$-linear part of $g_e$ is
\emph{not} the two-component operator; it is the two-component operator minus
the height-weighted connected operator $\Psi_{e,1}$.  The discrepancy
$\Psi_{e,1}$ is the merging term, and Theorem~\ref{thm:A} below shows it is the
term that carries all the operator content.
\end{remark}

\section{Two hook lemmas}

Both closed forms are computed on hooks, where the poset of subshapes is a
chain-like ladder and the enumeration is completely rigid.  We use throughout:
the coefficient of $s_\mu$ in $M_{p_1}^{\,r}\,X\,M_{p_1}^{\,d-r}s_\varnothing$ is
\begin{equation}\label{eq:F}
  F_X(r)\;=\;\sum_{\nu\subseteq\kappa\subseteq\mu}
     f^{\nu}\;\langle s_\kappa,\,X\,s_\nu\rangle\;f^{\mu/\kappa},
  \qquad |\nu|=d-r,\ \ |\mu/\kappa|=r,
\end{equation}
where $f^{\nu}$ (resp.\ $f^{\mu/\kappa}$) is the number of standard Young
tableaux of shape $\nu$ (resp.\ skew shape $\mu/\kappa$), because $M_{p_1}$ adds
one box with coefficient $1$.  Then
\begin{equation}\label{eq:diff}
  \langle s_\mu,\ \ad^{\,d}(X)\,s_\varnothing\rangle
  =\sum_{r=0}^{d}(-1)^r\binom{d}{r}F_X(r).
\end{equation}
The right side annihilates every polynomial in $r$ of degree $<d$.  Each of the
two lemmas below shows that $F_{R_e(t)}$ is such a polynomial \emph{plus a
spike at $r=d$}, and the spike is what survives.

Recall $f^{(a,1^b)}=\binom{a+b-1}{b}$, and that if a skew shape is a disjoint
union of two connected chains (single rows or single columns) of sizes $p$ and
$q$ then $f=\binom{p+q}{p}$.

\begin{lemma}[arm-$1$ hooks]\label{lem:hook2}
Let $\mu=(2,1^m)$.  Its subshapes are exactly $(1^j)$, $0\le j\le m+1$, and
$(2,1^j)$, $0\le j\le m$.  For $e\ge2$, the pairs $\nu\subseteq\kappa\subseteq\mu$
with $\kappa/\nu$ a \emph{connected} border strip of size $e$ are exactly:
\begin{enumerate}
\item[(A)] $\nu=(1^{j'})$, $\kappa=(1^{j'+e})$, with $\hgt=e-1$;
\item[(B)] $\nu=\varnothing$, $\kappa=(2,1^{e-2})$, with $\hgt=e-2$;
\item[(C)] $\nu=(2,1^{j'})$, $\kappa=(2,1^{j'+e})$, with $\hgt=e-1$;
\end{enumerate}
and those with $\kappa/\nu$ a broken border strip of size $e$ with exactly two
components are exactly
\begin{enumerate}
\item[(D)] $\nu=(1^{j'})$ with $j'\ge1$, $\kappa=(2,1^{j'+e-2})$, with $\hgt=e-2$.
\end{enumerate}
\end{lemma}

\begin{proof}
A subshape of $\mu$ has first part $\le2$ and all later parts $\le1$, which gives
the list.  If $\kappa=(1^j),\nu=(1^{j'})$ the skew shape is the column-$1$ cells
in rows $j'+1,\dots,j$: one component, $j-j'$ rows.  If
$\kappa=(2,1^j),\nu=(2,1^{j'})$ it is the column-$1$ cells in rows
$j'+2,\dots,j+1$: again one component.  If $\kappa=(2,1^j)$ and $\nu=(1^{j'})$
the skew shape is the single cell $(1,2)$ together with the column-$1$ cells in
rows $j'+1,\dots,j+1$.  For $j'=0$ this set contains $(1,1)$, which is
edge-adjacent to $(1,2)$: one component, $j+1$ rows, and size $j+2$; size $e$
forces $j=e-2$ and $\hgt=e-2$, which is (B).  For $j'\ge1$ the column-$1$ cells
all lie in rows $\ge2$ and are not edge-adjacent to $(1,2)$: two components, of
sizes $1$ and $j+1-j'$, with $\hgt=0+(j-j')$; size $e$ forces $j=j'+e-2$ and
$\hgt=e-2$, which is (D).  No shape here contains a $2\times2$ square, since
column $2$ carries at most one cell.  Sizes and heights in (A),(C) are as
stated.
\end{proof}

\begin{lemma}[hooks with arm $e$]\label{lem:hooke}
Let $\mu=(e,1^d)$ with $e\ge2$.  The pairs $\nu\subseteq\kappa\subseteq\mu$ with
$\kappa/\nu$ a connected border strip of size $e$ are exactly:
\begin{enumerate}
\item[(i)] $\nu=\varnothing$, $\kappa=(e-b,1^{b})$ for $0\le b\le\min(d,e-1)$,
with $\hgt=b$;
\item[(ii)] $\nu=(a,1^{b'})$ with $a\ge1$, $\kappa=(a,1^{b'+e})$, with
$\hgt=e-1$.
\end{enumerate}
\end{lemma}

\begin{proof}
Subshapes of $\mu$ are $\varnothing$ and $(a,1^b)$ with $1\le a\le e$,
$0\le b\le d$.  For $\nu=(a',1^{b'})$ and $\kappa=(a,1^b)$ the skew shape is the
row-$1$ cells in columns $a'+1,\dots,a$ together with the column-$1$ cells in
rows $b'+2,\dots,b+1$.  If $a'\ge1$ both blocks avoid the corner and are
mutually non-adjacent, so connectivity forces one of them to be empty; the
row-$1$ block alone has size $a-a'\le e-1<e$, so it is the column-$1$ block that
survives, giving (ii) with $b-b'=e$.  If $\nu=\varnothing$ the skew shape is all
of $\kappa=(a,1^b)$, a hook, which is connected, contains no $2\times2$, has
$b+1$ rows and size $a+b$; size $e$ gives (i) with $a=e-b\ge1$.
\end{proof}

\section{The two closed forms}

\begin{theorem}\label{thm:A}
For all $e\ge2$ and $d\ge1$, with $\mu=(2,1^{\,d+e-2})$,
\[
  \bigl\langle s_{\mu},\ \ad^{\,d}\bigl(R_e(t)\bigr)\,s_\varnothing\bigr\rangle
  \;=\;(-1)^{d}\,t^{\,e-2}\,(1+t).
\]
\end{theorem}

\begin{proof}
Write $m=d+e-2$, so $|\mu|=m+2=d+e$.  We evaluate \eqref{eq:F} for
$X=R_e(t)$ using Lemma~\ref{lem:hook2}(A)--(C).  In each case
$r=|\mu|-|\kappa|$ and $|\nu|=d-r$.

\emph{(A)} $\nu=(1^{j'})$, $\kappa=(1^{j'+e})$, so $r=(m+2)-(j'+e)=d-j'$, i.e.\
$j'=d-r$.  The constraint $\kappa\subseteq\mu$ reads $j'+e\le m+1=d+e-1$, i.e.\
$j'\le d-1$, i.e.\ $r\ge1$.  Here $f^{\nu}=1$; and $\mu/\kappa$ consists of the
single cell $(1,2)$ together with $m+1-(j'+e)$ cells in column $1$ starting at
row $j'+e+1\ge3$, two disjoint chains, so
$f^{\mu/\kappa}=\binom{1+(m+1-j'-e)}{1}=m+2-j'-e=d-j'=r$.  The weight is
$t^{e-1}$.  Total: $F_A(r)=r\,t^{e-1}$ for $1\le r\le d$, and $F_A(0)=0$, so
$F_A(r)=r\,t^{e-1}$ for all $0\le r\le d$.

\emph{(B)} $\nu=\varnothing$, $\kappa=(2,1^{e-2})$, $|\kappa|=e$, so $r=d$ and
$|\nu|=0$, consistent.  Here $f^{\nu}=1$, $\mu/\kappa$ is the column-$1$ chain in
rows $e,\dots,m+1$ so $f^{\mu/\kappa}=1$, and the weight is $t^{e-2}$.  Thus
$F_B(r)=t^{e-2}\,\ind{r=d}$.

\emph{(C)} $\nu=(2,1^{j'})$, $\kappa=(2,1^{j'+e})$, so $r=(m+2)-(j'+e+2)=d-2-j'$.
The constraint $\kappa\subseteq\mu$ reads $j'+e\le m$, i.e.\ $j'\le d-2$, i.e.\
$r\ge0$; and $j'\ge0$ gives $r\le d-2$.  Here $f^{\nu}=f^{(2,1^{j'})}=j'+1=d-1-r$,
$\mu/\kappa$ is a single column chain so $f^{\mu/\kappa}=1$, and the weight is
$t^{e-1}$.  Thus $F_C(r)=(d-1-r)\,t^{e-1}$ for $0\le r\le d-2$ and $0$ otherwise.

Adding, for $r\le d-2$ we get $\bigl(r+(d-1-r)\bigr)t^{e-1}=(d-1)t^{e-1}$; for
$r=d-1$ we get $(d-1)t^{e-1}$; and for $r=d$ we get
$d\,t^{e-1}+t^{e-2}=(d-1)t^{e-1}+t^{e-2}(1+t)$.  Hence
\[
  F_{R_e}(r)=(d-1)\,t^{e-1}\;+\;t^{e-2}(1+t)\,\ind{r=d}.
\]
Now apply \eqref{eq:diff}: the constant term is annihilated because
$\sum_{r}(-1)^r\binom dr=0$ for $d\ge1$, and the spike contributes
$(-1)^d\binom dd\,t^{e-2}(1+t)$.
\end{proof}

\begin{theorem}\label{thm:B}
For all $e\ge2$ and $d\ge1$, with $\mu=(e,1^{d})$,
\[
  \bigl\langle s_{\mu},\ \ad^{\,d}\bigl(R_e(t)\bigr)\,s_\varnothing\bigr\rangle
  \;=\;(-1)^{d}\Biggl[\ \sum_{b=0}^{\min(d,\,e-1)}\binom{d}{b}t^{b}
     \;-\;\binom{d-1}{e-1}t^{\,e-1}\Biggr].
\]
\end{theorem}

\begin{proof}
Here $|\mu|=e+d$.  Use Lemma~\ref{lem:hooke}.

\emph{(i)} $\nu=\varnothing$, $\kappa=(e-b,1^{b})$, $|\kappa|=e$, hence $r=d$ and
$|\nu|=0$.  Then $f^{\nu}=1$ and $\mu/\kappa$ consists of the $b$ row-$1$ cells in
columns $e-b+1,\dots,e$ and the $d-b$ column-$1$ cells in rows $b+2,\dots,d+1$,
disjoint chains, so $f^{\mu/\kappa}=\binom{d}{b}$.  The weight is $t^{b}$.  Thus
case (i) contributes $\ind{r=d}\sum_{b=0}^{\min(d,e-1)}\binom{d}{b}t^{b}$.

\emph{(ii)} $\nu=(a,1^{b'})$, $\kappa=(a,1^{b'+e})$, so
$r=(e+d)-(a+b'+e)=d-a-b'$; write $s:=a+b'=d-r$.  Then
$f^{\nu}=\binom{s-1}{b'}=\binom{s-1}{a-1}$, and $\mu/\kappa$ consists of the
$e-a$ row-$1$ cells in columns $a+1,\dots,e$ and the $d-b'-e$ column-$1$ cells in
rows $b'+e+2,\dots,d+1$, so $f^{\mu/\kappa}=\binom{r}{e-a}$ --- an expression
that vanishes automatically when $a>e$ or when $b'+e>d$, so no further
constraints are needed.  The weight is $t^{e-1}$.  For $s\ge1$, Vandermonde gives
\[
  \sum_{a\ge1}\binom{s-1}{a-1}\binom{r}{e-a}
  =\sum_{i\ge0}\binom{s-1}{i}\binom{r}{e-1-i}
  =\binom{s-1+r}{e-1}=\binom{d-1}{e-1},
\]
a constant independent of $r$.  For $s=0$, i.e.\ $r=d$, the sum is empty and
equals $0$.  Hence
\[
  F_{R_e}(r)=\binom{d-1}{e-1}t^{e-1}
  \;+\;\ind{r=d}\Biggl[\sum_{b=0}^{\min(d,e-1)}\binom{d}{b}t^{b}
        -\binom{d-1}{e-1}t^{e-1}\Biggr],
\]
and \eqref{eq:diff} kills the constant and keeps $(-1)^d$ times the bracket.
\end{proof}

\begin{remark}
The two theorems agree where the two hook families coincide.  At $e=2$ both
$\mu$'s are $(2,1^{d})$, and Theorem~\ref{thm:B} gives
$(-1)^d\bigl[1+dt-(d-1)t\bigr]=(-1)^d(1+t)$, which is Theorem~\ref{thm:A} at
$e=2$.
\end{remark}

\begin{remark}[the mechanism]\label{rem:mech}
In both proofs $F_{R_e}(r)$ is a polynomial in $r$ of degree $<d$ plus a spike
at $r=d$, and the spike is the contribution of the chains with $\nu=\varnothing$:
those in which the whole $e$-ribbon is added to the \emph{empty} partition.
This is exactly the locus where ``connected'' and ``broken'' disagree, because
the corner cell $(1,1)$ glues the arm to the leg.  The alternating sum
$\sum_r(-1)^r\binom dr$ is blind to polynomials and sees only the spike, so
\emph{the whole theorem is one exceptional chain}.  In Theorem~\ref{thm:A} the
spike is $t^{e-2}+t^{e-1}$: the ribbon added to $\varnothing$ can sit as the hook
$(2,1^{e-2})$ (height $e-2$) or, as the failure of case (C) at $r=d$, as the
column $(1^{e})$ (height $e-1$).  \emph{The factor $(1+t)$ is those two seats.}
\end{remark}

\section{Consequences}

\begin{corollary}\label{cor:ordR}
Let $c$ be a scalar.  Then $\ord\bigl(R_e(c)\bigr)=0$ if $c=-1$, and
$\ord\bigl(R_e(c)\bigr)=\infty$ for every $c\neq-1$.  In particular
$\ord_{\Q(t)}\bigl(R_e(t)\bigr)=\infty$ for every $e\ge2$.
\end{corollary}

\begin{proof}
$R_e(-1)=M_{p_e}$ has order $0$.  For $c\notin\{0,-1\}$, Theorem~\ref{thm:A}
gives $\ad^{\,d}(R_e(c))\neq0$ for every $d\ge1$, since
$(-1)^dc^{e-2}(1+c)\neq0$.  For $c=0$, Theorem~\ref{thm:B} gives
$(-1)^d\bigl[\binom d0\bigr]=(-1)^d\neq0$ (all other terms carry a positive
power of $t$, and $e-1\ge1$).  Apply Lemma~\ref{lem:gen}.
\end{proof}

\begin{corollary}\label{cor:ordN}
For every $e\ge2$, $\ord(\Psi_{e,1})=\ord(N_e)=\ord(E_e)=\infty$.  Explicitly,
for all $d\ge1$,
\[
  \bigl\langle s_{(2,1^{d+e-2})},\ \ad^{\,d}(N_e)\,s_\varnothing\bigr\rangle
  =(-1)^{\,d+e+1}.
\]
In particular $N_e$ is not a derivation, is not of finite order, and is not a
differential operator on $\LL$.
\end{corollary}

\begin{proof}
The operators $M_{p_1}$ do not depend on $t$, so $\ad$ commutes with
$\partial_t$.  Applying $-\partial_t|_{t=-1}$ to Theorem~\ref{thm:A} and using
$\Psi_{e,1}=-\partial_tR_e|_{t=-1}$ gives
\[
  \bigl\langle s_{(2,1^{d+e-2})},\ \ad^{\,d}(\Psi_{e,1})\,s_\varnothing\bigr\rangle
  =-\,(-1)^{d}\,\partial_t\bigl(t^{e-2}(1+t)\bigr)\Big|_{t=-1}
  =-\,(-1)^{d}\,(-1)^{e-2}=(-1)^{\,d+e+1},
\]
where the middle evaluation holds because $(1+t)$ vanishes at $t=-1$, so only
the term in which $\partial_t$ hits the factor $(1+t)$ survives, leaving
$t^{e-2}|_{t=-1}$.
By Proposition~\ref{prop:merge}, $\ad^{\,d}(N_e)=\ad^{\,d}(\Psi_{e,1})$ for
$d\ge1$, which gives the displayed value; it is never zero, so
$\ord(N_e)=\infty$ by Lemma~\ref{lem:gen}.  A derivation has order $\le1$, so
$N_e$ is not one.  Finally $E_e=\eps N_e+O(\eps^2)$ and $\ord$ of every
$t$-Taylor coefficient of $E_e$ is bounded by $\ord_{\Q(t)}(E_e)$, so
$\ord(E_e)=\infty$; alternatively $E_e=g_e-R_e$ with $\ord(g_e)=0$ and
Corollary~\ref{cor:ordR}.
\end{proof}

\begin{corollary}[no Pieri presentation]\label{cor:pieri}
Let $\mathcal A\subseteq\operatorname{End}_K(\LL)$ be the (associative, unital)
algebra generated by all multiplication operators $M_f$ and all skewing
operators $f^{\perp}$, $f\in\LL$.  Then for every $e\ge2$ and every scalar
$c\neq-1$,
\[
  R_e(c)\notin\mathcal A,\qquad E_e(c)\notin\mathcal A,\qquad N_e\notin\mathcal A .
\]
\end{corollary}

\begin{proof}
By Lemma~\ref{lem:std}, $\ord(M_f)=0$ and $\ord(f^{\perp})=\deg_pf<\infty$, and
order is subadditive under composition and takes the max under sums; so every
element of $\mathcal A$ has finite order.  Corollaries~\ref{cor:ordR} and
\ref{cor:ordN} say these operators do not.
\end{proof}

\begin{remark}[what this settles for the $(1+t)$ programme]
\begin{enumerate}
\item The route to the sharpness of $(1+t)^{k-1}$ through ``is $N_e$ a
derivation?'' is closed, and closed in the strongest direction: $N_e$ is not
merely a non-derivation, it has no finite order at all.  The predicate
$j_{\min}\ge\lceil (k-1)/d\rceil$ with $d=\ord N_e$ degenerates to
$j_{\min}\ge1$, which is consistent with --- and gives no new information
beyond --- the independently proved sharpness statement.
\item $\ord(g_e)=0$ while $\ord(R_e)=\ord(E_e)=\infty$: the decomposition of a
multiplication operator into its connected and disconnected parts destroys
finite order.  This is a quantitative form of ``a truncation is not a
multiplication operator'', and it explains structurally why the commutator of
two such truncations is not addressed anywhere in the Hall--Littlewood
literature: the objects are not differential operators, so none of the standard
machinery for $M_f$ and $f^{\perp}$ applies to them.
\item Corollary~\ref{cor:pieri} \emph{separates} the ribbon operators from the
algebra generated by Pieri operators and their transposes, rather than reducing
one to the other.  Any question posed inside that algebra --- e.g.\ the
commutation relations between $h_i\cdot$ and $h_j^{\perp}$ --- lives entirely in
finite order and therefore cannot have $R_e$ or $N_e$ among its consequences.
\item The order filtration of $R_e(t)$ is \emph{discontinuous at the anchor}:
$0$ at $t=-1$, $\infty$ everywhere else.  Every measurement made at $t=-1$ is
therefore blind to the operator content of the family.
\end{enumerate}
\end{remark}

\begin{remark}[the bilinearity identity, for the record]\label{rem:q87}
Since $g_e=R_e+E_e$ and $[g_e,g_{e'}]=0$ (multiplication operators in a
commutative ring, Proposition~\ref{prop:hl}), bilinearity of the bracket gives
\[
  [R_e,R_{e'}]\;=\;-\bigl([R_e,E_{e'}]+[E_e,R_{e'}]+[E_e,E_{e'}]\bigr).
\]
This is one line and it cannot fail; it carries no information beyond
$[g_e,g_{e'}]=0$, which is itself forced.  Its only genuine consequence is the
$\eps^0$ statement: both sides vanish at $t=-1$ because $E_e(-1)=0$ and
$R_e(-1)=M_{p_e}$, which re-derives $(1+t)\mid[R_e,R_{e'}]$ in one step.  It is
recorded here so that it is not mistaken for content.
\end{remark}

\begin{remark}[a reduction whose two sides are both infinite]\label{rem:vac}
A previous session recorded, from $\ord(g_e)=0$, that
$\ord(R_e)=\ord(E_e)$ and concluded that ``the operator content lives entirely
in the disconnected part''.  The identity is correct and is reproved by
Corollary~\ref{cor:ordN}.  But both sides are now known to be $\infty$, so the
identity bounds nothing: it is a true equation between two infinities, not a
reduction.  What survives of that observation is the sharper
Corollary~\ref{cor:pieri}.
\end{remark}

\section{Verification}

All computations are in \verb|probes/2026-09-06-Q84/|.  Two engines were used and
kept code-disjoint: \verb|engine.py| enumerates skew shapes directly and works
over $\Q(t)$ with \textsf{sympy}; \verb|fast.py| works at $t=-1$ (or at integer
$t$) with the Maya/abacus bead model for connected ribbons and a constructive
row-increment enumeration for broken ones.

\begin{enumerate}
\item \emph{$R_e$ against prior code.}  \verb|engine.op_R| agrees with
\verb|probes/2026-09-04-Q76/abacus.py| on all $225$ pairs
$(\lambda,e)$ with $|\lambda|\le7$, $1\le e\le5$.
\item \emph{Proposition~\ref{prop:hl}.}  $M[P_{(e)}(X;-t)]s_\lambda=g_e s_\lambda$
on all $150$ pairs with $1\le e\le5$, $|\lambda|\le6$, exactly, with no free
scalar: the left side is computed from \eqref{eq:Prow} by iterated
Murnaghan--Nakayama, the right side from broken-strip enumeration.
\item \emph{Proposition~\ref{prop:merge}.}  $N_e-\Psi_{e,1}=M[\cdots]$ on all
$335$ pairs with $2\le e\le6$, $|\lambda|\le8$.
\item \emph{$N_e s_\varnothing=0$}, checked (not assumed) for $e=3,4,5$.
\item \emph{Theorem~\ref{thm:A}.}  $245/245$ agreements over
$e\in\{2,\dots,6\}$, $d\in\{1,\dots,7\}$, $t\in\{-3,-2,-1,0,1,2,4\}$; and the
symbolic version over $\Q(t)$ for $e\le4$, $d\le3$.
\item \emph{Corollary~\ref{cor:ordN}.}  $(-1)^{d+e+1}$ confirmed on $40/40$
pairs $e\in\{2,\dots,6\}$, $d\in\{1,\dots,8\}$, computed from the
two-component definition \eqref{eq:Ndef} directly --- a code path disjoint from
the one used for Theorem~\ref{thm:A}.
\item \emph{Theorem~\ref{thm:B}.}  $280/280$ agreements over
$e\in\{2,\dots,6\}$, $d\in\{1,\dots,7\}$, $t\in\{-3,\dots,5\}$.
\item \emph{Negative controls} (\verb|controls.py|).  The same bracket harness,
run on operators of known order, vanishes at exactly the right depth:
$M_{h_2}$ (order $0$) has $C^{(1)}=0$; $M_{p_1}p_1^{\perp}$ (order $1$) has
$C^{(1)}\neq0$, $C^{(2)}=0$; $M_{p_3p_1p_1}(p_1^{\perp})^2$ (order $2$, degree
$+3$, structurally parallel to $N_3$) has $C^{(2)}\neq0$, $C^{(3)}=0$.  The
sharpest control is the minimal pair: $M_{h_2}$ and $N_2$ are computed by the
same two-box enumeration with one predicate flipped (horizontal strip versus two
components), and have orders $0$ and $\infty$.
\item \emph{Full sweep.}  $C^{(d)}\neq0$ for $N_e$, $e\in\{2,3,4\}$,
$d\in\{1,2,3\}$, over all ordered $a$-tuples in $\{1,2,3\}^d$
(\texttt{itertools.product}, not \texttt{combinations}) and all $\lambda$ with
$|\lambda|\le6$: $3481$ of $3510$ triples $(e,\mathbf a,\lambda)$ give a nonzero
result.  The $29$ vanishing triples are individual $(\lambda,\mathbf a)$
coincidences spread over $|\lambda|\in\{1,3,5,6\}$, not a pattern; what is at
stake is the vanishing of an \emph{operator}, and Theorems~\ref{thm:A} and
\ref{thm:B} exhibit states where it does not vanish, for every $d$.
\end{enumerate}

\section{Gaps and what is not claimed}

\begin{enumerate}
\item Corollary~\ref{cor:ordR} is stated for scalars $c$ in a field of
characteristic $0$; the two hook families jointly cover $c\neq-1$, but a single
uniform witness family covering all $c\neq-1$ at once is not given.
\item Nothing here computes $\ord(\Psi_{e,D})$ for $D\ge2$.  The natural
guess --- $0$ for $D=0$ and $\infty$ for every $D\ge1$ --- is untested.
\item Proposition~\ref{prop:hl} is used only through $\ord(g_e)=0$ and the
$\eps$-linear coefficient of \eqref{eq:Prow}; both are verified independently
above.  The attribution of $g_e=M[P_{(e)}(X;-t)]$ to the quantum
Murnaghan--Nakayama rule rests on a prior reading session, but the identity
itself is derived here from the Hall--Littlewood generating function and
checked, so no step of the present argument depends on the attribution.
\item Remark~\ref{rem:mech}(3) refers to the algebra generated by Pieri
operators and their transposes.  What is proved is the mathematical statement
Corollary~\ref{cor:pieri}; the reading of it as an answer to a particular posed
question is a remark, not a theorem, and the source of that question was not
re-read in this session.
\end{enumerate}

\end{document}
